\pdfoutput=1
\documentclass[11pt]{article}
\usepackage[T1]{fontenc}
\usepackage{lmodern}
\usepackage[margin=1in]{geometry}
\usepackage{amsmath,amssymb,amsthm,mathtools,microtype}
\usepackage[hidelinks]{hyperref}
\hypersetup{pdftitle={A polarization symplectic form at a fixed two-stream null-dust limit}}
\newtheorem{theorem}{Theorem}
\newtheorem{lemma}[theorem]{Lemma}
\newcommand{\tr}{\operatorname{tr}}
\newcommand{\re}{\operatorname{Re}}
\newcommand{\im}{\operatorname{Im}}
\newcommand{\dd}{\mathrm d}
\newcommand{\supp}{\operatorname{supp}}
\title{A polarization symplectic form at a fixed two-stream null-dust limit\\
\large An exact characteristic high-frequency construction}
\author{}
\date{Paper 29C}
\begin{document}
\maketitle

\begin{abstract}
We construct two independently variable gravitational-wave
polarizations in exact vacuum developments with a single
double-null weak limit and the same two resolved scalar
null-dust measures. The two streams can overlap in the
interior. On the marked characteristic Cauchy surface the
unscaled Einstein Lee--Wald form tends to zero on the
selected bounded profile tangents, but division by the
reduced wavelength leaves two independent weighted projective
polarization forms. A slow amplitude correction cancels an
otherwise nonzero first-order focusing term in the exact
null constraint. This gives the differentiated area
estimates used in the proof of the symplectic limit. The
limiting form is intrinsic to the specified rank-one
polarization data on the two initial sheets and transports
to a spacelike cut with the same edges. A scalar-dust
description omits these pairings.
\end{abstract}

\section{Statement and relation to prior work}

Use intersecting null sheets \(H_1=\{u=0,\,0\leq v\leq L_1\}\)
and \(H_2=\{v=0,\,0\leq u\leq L_2\}\), with spherical
sections and the same marked initial sphere. On each sheet
\(s_i\) is its own fixed increasing affine parameter
(\(v\) on \(H_1\), \(u\) on \(H_2\)). The generator labels,
the outer-cut embeddings and the incident lapse are fixed.
We take the incident lapse equal to one. All spacetime
comparisons use the same double-null identification; no
wavelength-dependent relabeling enters a variation.

On \(H_i\) fix a smooth positive slow screen metric
\(\sigma_i\), with \(\partial_{s_i}\dd A_{\sigma_i}=0\),
and a real smooth envelope \(F_i\) supported in an angular
frame patch and away from both generator endpoints.
The slow data agree on the common sphere, and the
corner expansions and torsion are fixed. The polarized
profiles \(z_i\) are smooth maps to the unit vectors in
\(\mathbb C^2\) in a fixed smooth orthonormal frame of real
tracefree symmetric screen endomorphisms. Only their
values on the envelope support are used. We take bounded
smooth profile charts, allowing arbitrary localized
variations in \(s_i\) and angle and compact finite-dimensional
families of such variations. At each wavelength a tangent
means a derivative of the exact polarized families, not a
substitution of an unconstrained test strain.

\begin{theorem}[Two-stream polarization pairing]\label{thm:main}
There are fixed data of the stated kind and a common
regular double-null rectangle for which, for every
sufficiently small reduced wavelength \(\lambda>0\) and profiles
in the chosen bounded charts, the following hold.
\begin{enumerate}
\item One obtains smooth exact vacuum metrics \(g_\lambda(z_1,z_2)\).
They have a common double-null limit, in \(C^0\) and weak
\(W^{1,2}\) on the common rectangle, independent of both
profiles. The two resolved scalar null-dust measure
families in that limit are independent of both profiles.
The data can be chosen so that both dusts have positive
density on a common open region.
\item Let \(\mathfrak S_\lambda\) be the bare Einstein
Lee--Wald two-form on the pair of initial sheets, with the
future Cauchy orientation specified in
Section~\ref{sec:current}. Write \(a_{0,i}^2\sigma_i\)
for the limiting induced metric there. On the selected
profile tangents,
\begin{equation}\label{eq:theorem}
 \lim_{\lambda\downarrow0}\lambda^{-1}\mathfrak S_\lambda
 =\sum_{i=1}^2\int_{H_i}
     \frac{a_{0,i}^2 F_i^2}{64\pi G}\,
     \dd_{\!P}\im(z_i^\dagger\dd_{\!P}z_i)\,
     \dd s_i\dd A_{\sigma_i}.
\end{equation}
Here \(\dd_{\!P}\) is exterior differentiation on profile
space. The unscaled form converges to zero on these tangents.
\item On \(F_i\ne0\), the two integrands descend to the
projective lines \([z_i]\in\mathbb{CP}^1\); with
\(z_i=e^{i\chi_i}(\sqrt{1-p_i},e^{i\phi_i}\sqrt{p_i})\)
the projective two-form is
\(\dd_{\!P}p_i\wedge\dd_{\!P}\phi_i\).
The two blocks have nonzero pairings for compactly
supported projective variations separately, with zero
cross pairing between the incident sheets. Pullback of the
same Lee--Wald current to an interior spacelike surface
with the same two outer cuts gives the same formula.
\end{enumerate}
\end{theorem}

The geometric data and the bounded family are chosen
on the short, noncaustic branch of the characteristic
existence theorem. The proof does not use a symmetry
reduction. The formula is local where the density is
positive, and it is not a prescription for a pointwise
bulk Jones polarization field. The normalization of
the projective form in \eqref{eq:theorem} has total
area \(2\pi\) on \(\mathbb{CP}^1\).

\paragraph{Prior results and the new comparison.}
Luk and Rodnianski construct and characterize
double-null high-frequency limits with two
measure-valued dusts \cite{LR}. Their exact constraints,
uniform existence, compactness, transport and uniqueness
are analytic inputs here. In the associated public Paper
29B \cite{paper29B}, exact two-stream data and directed
Hermitian polarization measures at fixed scalar trace
already appear, including a broader smooth factorizable
matrix class and an optical readout. We use a single
coherent carrier per incident face. This paper adds an
Einstein symplectic comparison on actual constrained
tangents, rather than another claim of two-stream
realizability. The complementary marked-cut result
was already mentioned without a proof in Paper 29
\cite{paper29}; the standalone derivation is given here. The relation to the finite-band
metaplectic probe transfer of \cite{transfer} is likewise
not an identification: that work concerns a probe on an
exact sandwich, whereas the form here compares
derivatives of the rapidly oscillating vacuum background.
The null gravitational symplectic form is prior work
\cite{Reisenberger,HF}; the calculation below fixes its
normalization and edge convention for the data used here.
The basic projective two-form is the standard one for
unit complex rays \cite{AS}; the new result is its
calibrated occurrence as a degenerate high-frequency
limit of the Einstein form at fixed effective geometry.

\section{Exact vacuum characteristic data}\label{sec:data}

We first work on one sheet and suppress its index.
Fix \(\sigma\) and the envelope \(F\).
Let \(E_1,E_2\) be the chosen real orthonormal tracefree
screen endomorphisms, and put
\[
 Z=z^1 E_1+z^2 E_2,\qquad \tr(Z^\dagger Z)=1,\qquad
 A=\sigma^{-1}\sigma'.
\]
A prime is the derivative in the stated generator
coordinate. The parameter \(\lambda\) is inverse affine
angular frequency: one phase period has length \(2\pi\lambda\). The metric-compatible derivative on
endomorphisms is \(D_s Z=Z'+[A,Z]/2\) in fixed
angular coordinates. Define
\begin{align}
 \beta&=\im\tr(Z^\dagger D_s Z),\nonumber\\
 H_\lambda&=\lambda F(1-\lambda\beta)
                    \re(Z e^{is/\lambda}),\qquad
 \widehat q_\lambda=\sigma e^{H_\lambda},\qquad
 q_\lambda=a_\lambda^2\widehat q_\lambda .
 \label{eq:ansatz}
\end{align}
The expression is extended as zero strain outside the
angular frame patch. The real strain is
\(\sigma\)-self-adjoint and tracefree, so the conformal
metric is positive and has the stated fixed determinant.
The compensating factor is a slow real change of
amplitude at order \(\lambda\), not a change of carrier
phase. The overall intensity to leading order is fixed.
We may take \(\lambda\) so small that this factor is
positive throughout the chosen profile chart.

The reduced vacuum affine Raychaudhuri equation is
\begin{equation}
 B_\lambda=\widehat q_\lambda^{-1}\widehat q_\lambda',
 \qquad
 a_\lambda''+Q_\lambda a_\lambda=0,\qquad
 Q_\lambda=\tfrac18\tr(B_\lambda^2).
 \label{eq:ray}
\end{equation}
It is precisely \cite[Eqs.~(3.1)--(3.2)]{LR} with incident
lapse one. The positive initial area and its initial slope
are fixed at the corner. This equation is solved
exactly, with no replacement of the shear by a weak
limit in the vacuum data.

\subsection{Canceling the first-order focusing drift}

At fixed slow coordinates write \(\vartheta=s/\lambda\),
\(U=F\re(Ze^{i\vartheta})\), and
\(V=F\re(iZe^{i\vartheta})=\partial_\vartheta U\).
Let \(U'_0\) denote a slow derivative holding
\(\vartheta\) fixed.
Expanding
\(e^{-H_\lambda}A e^{H_\lambda}
 +e^{-H_\lambda}(e^{H_\lambda})'\)
in~\eqref{eq:ray} gives
\begin{equation}
 B_\lambda=A+V+
   \lambda\bigl(U'_0+[A,U]+\tfrac12[V,U]-\beta V\bigr)
   +O(\lambda^2).
 \label{eq:B}
\end{equation}
The expansion and the error bound hold after any fixed
angular number and two profile derivatives on a fixed
bounded smooth chart. In particular one must retain the
\(-\beta V\) term from differentiating the corrected
strain. At fixed \(\sigma\), each of \(A,V,U\) is
self-adjoint (the last two are real), so the commutators
in \eqref{eq:B} are skew-adjoint and do not contribute
at order \(\lambda\) to \(\tr(B_\lambda^2)\). Replacing
\(U'_0\) by \(D_s U\) in that contraction is legitimate
for the same reason.

Let angled brackets mean averaging over \(\vartheta\)
only, at fixed slow \(s,\theta\) and profile. Then
\[
 \langle \tr(V^2)\rangle=\tfrac12 F^2,\qquad
 \langle\tr(V D_sU)\rangle
  =\tfrac12 F^2\im\tr(Z^\dagger D_s Z)
  =\tfrac12F^2\beta .
\]
The \(F'\) term in the second average has zero real
contraction with \(iZ\). The mean from \(V D_sU\)
would add \(\lambda F^2\beta/8\) to \(Q_\lambda\)
without the compensator. The term \(-\lambda\beta V\)
cancels it exactly at this order. Therefore
\begin{equation}
 Q_\lambda-Q_0
 =\sum_{0<|m|\leq 2} (q_m+\lambda\widetilde q_m)
                         e^{ims/\lambda}+O(\lambda^2),
 \qquad Q_0=\tfrac18\tr(A^2)+\tfrac1{16}F^2 .
 \label{eq:average}
\end{equation}
The harmonic coefficients are slow, smooth and supported
away from the generator endpoints. The same bounds as in
\eqref{eq:B} hold for the error and the specified
differentiated coefficients. In particular the mean
\(Q_0\) and the initial data do not depend on either
polarization.

\begin{lemma}[Differentiated area estimate]\label{lem:area}
Let \(a_0''+Q_0a_0=0\) with the same initial conditions.
On a fixed finite interval where \(a_0\) is positive,
\begin{equation}
   a_\lambda-a_0=O(\lambda^2),\qquad
   a'_\lambda-a'_0=O(\lambda).
 \label{eq:area}
\end{equation}
The estimates hold through two profile derivatives and
any fixed angular order. Beyond the generator support of
\(F\), the slope estimate is \(O(\lambda^2)\).
\end{lemma}

\begin{proof}
Write the difference as the retarded Volterra integral
with source \(-(Q_\lambda-Q_0)a_0\) and kernel
\(G_\lambda(s,t)\) of \(\partial_s^2+Q_\lambda\).
Construction from two homogeneous solutions and
their constant Wronskian gives
\(G_\lambda(s,s)=0\), \(\partial_sG_\lambda(s,s)=1\),
and, for \(t<s\),
\(\partial_t^2G_\lambda=-Q_\lambda(t)G_\lambda\).
The last identity holds for \(\partial_sG_\lambda\)
as well. Since \(Q_\lambda\) and the stipulated
angular and profile derivatives are bounded, direct
first-order ODE estimates bound the kernels and
their first two \(t\)-derivatives, with those same
parameter derivatives. This argument does not
differentiate the rapid source in \(t\).
For each nonzero harmonic in \eqref{eq:average},
integrate by parts twice in \(t\) for the area;
the first boundary term vanishes by the corner
buffer and \(G_\lambda(s,s)=0\). For the slope one
integration is always sufficient, and its first
boundary term also vanishes beyond the envelope
support. The remainder in \eqref{eq:average}
has \(O(\lambda^2)\) contribution even without
integration by parts. Differentiate the integral
with respect to the smooth parameters and angles:
the oscillatory phase is unchanged and all the
required kernel identities remain bounded by the
same ODE argument. This proves the lemma.
\end{proof}

Consequently the exact \(a_\lambda\) stay uniformly
positive on a short common interval, and the variation
obeys \(\delta a_\lambda=O(\lambda^2)\),
\(\delta a_\lambda'=O(\lambda)\).
These are differentiated estimates for the exact
constraint solution, not an inference merely from
strong convergence of \(q_\lambda\).

\subsection{The remaining constraints and the common limit}

We use the double-null gauge of \cite[Eq.~(2.2) and
Def.~2.10]{LR}, with \(e_3=\Omega^{-1}(\partial_u+b)\)
and \(e_4=\Omega^{-1}\partial_v\). In the following description
all equation numbers except ours refer to \cite{LR}.
On \(H_1\), the intrinsic \(q_\lambda\) fixes the outgoing
expansion and shear. Solve the torsion equation (2.15) with
its fixed corner value, recovering the conjugate torsion
from their prescribed sum \(2\nabla\log\Omega=0\). Then
reconstruct \(b\) by (2.9), solve the opposite expansion
and shear using (2.17) and (2.19), and finally use (2.21)
for the opposite lapse coefficient. On \(H_2\), \(b=0\)
by the initial gauge; use (2.16), (2.18), (2.20), and
finally (2.22) in the corresponding order. The last
equation uses the already reconstructed mixed shear.
The incident lapse is identically one on each initial
face, so its generator is affine. The opposite quantities
start at their fixed corner values, not at a separately
varying slice of the other wave. The reduced data agree
with the fixed slow data on an entire corner buffer; the
joint is exactly compatible.

Here is the convergence check for these additional
traces. The intrinsic metrics and generator expansions
converge strongly with every fixed reserved angular
derivative. The generator shear is the common slow
shear plus a bounded first harmonic and \(O(\lambda)\);
an angular derivative preserves that description. In the
subsequent triangular equations the incident fast shear
and its angular derivatives occur linearly. They can
multiply opposite-shear or other preceding transport
quantities; the Raychaudhuri term already treated is
the only same-face fast shear square. In fixed angular
components the connection along a generator also
contains only first metric derivatives, and therefore
has this same decomposition. At a transport step write
the difference from the common averaged transport
solution in integral form
\[
 e_\lambda(s)=\int_0^s C_\lambda(t)e_\lambda(t)\,\dd t+
 R_\lambda(s).
\]
The coefficients \(C_\lambda\) are uniformly bounded.
Terms from previously solved differences are
\(O(\lambda)\), and terms linear in the fast harmonics
against the common smooth slow solution have
\(O(\lambda)\) primitives by one integration by parts.
Thus \(R_\lambda=O(\lambda)\) uniformly, and Gronwall
closes the step. Angular differentiation gives the same
estimates with finitely many derivatives reserved for the
curvature, divergence and gradient terms. This argument
gives the common boundary traces, including the shift
and both opposite expansions and the incoming torsion
difference, not just the common induced metric.

Choose a fixed smooth corner stock of positive radius
and bounded angular geometry, and take the strips short
enough for all areas and the limiting metric to stay
regular. Lemma 3.1 and Theorems 3.2--3.3 of \cite{LR}
give exact smooth vacuum developments on one common
rectangle: the initial angular norms and generator
\(L^2\) shear bounds are uniform; second generator
derivatives of the conformal data need not be uniform.
We place the packets in that common short existence
window. Concretely, fix angular norms and amplitude
bounds on broader initial strips, use the uniform
existence estimates to choose a common short output,
and put the envelope supports in its interior. Shrinking
a chosen support does not increase the leading shear
bound. Terms involving its slow generator derivatives
are multiplied by positive powers of \(\lambda\) in
the conformal derivative, and are bounded by decreasing
the wavelength cutoff. This does not assert that
arbitrarily long prescribed envelopes are reached by the
short-time theorem. At fixed \(\lambda\), smooth
finite-dimensional profile variations yield actual
linearized vacuum tangents by smooth characteristic
dependence (or equivalently by its differentiated
characteristic equations). No uniform differentiated
interior convergence theorem is being assumed.

In the affine normalization the incident shear square is
\(\tr(B_\lambda^2)/4\). The cross term with the common
slow shear tends weakly to zero and the squared fast
part has mean \(F_i^2/2\). Thus the additional measure on
each incident sheet is exactly
\begin{equation}\label{eq:dust}
 \nu_i^{\rm in}
 =\frac{F_i^2}{8}a_{0,i}^2\,\dd s_i\,\dd A_{\sigma_i}.
\end{equation}
This is the normalization for
\(\Omega^2|\widehat\chi|^2\,\dd s\,\dd A_q\) (and its
ingoing counterpart) in \cite{LR}. The smooth slow
shear square is not part of the additional measure.
The explicit initial square is uniformly bounded and
does not concentrate at the initial corners. Use the
two-sheet compactness theorem (Theorem 4.1 of \cite{LR})
with its auxiliary transport equations, and the initial
trace comparison of its Theorem 4.4 and Section 9.4
(in particular the strong/weak trace inputs in (9.61)--(9.62)):
the strong initial transports just established supply the
same opposite expansion and torsion traces and
\eqref{eq:dust} supplies the two identical measure
traces. No trace of a differentiated interior weak limit
is inferred by formally differentiating its induced
metric. Uniqueness in that class makes the metric and
both complete measure families common to all
subsequential limits. Compactness and this uniqueness
give the claimed \(C^0\) and weak \(W^{1,2}\)
convergences of the family.

For actual two-stream overlap, use
\cite[Prop.~7.7]{LR}: the measure from \(H_1\) is
transported by the angular flow of the limiting shift
at fixed \(v\), and the measure from \(H_2\) by the
unchanged angular labels at fixed \(u\). Pick intervals
in the common rectangle where each envelope is positive
on a common angular patch, with a smaller patch separated
from its edge by more than the flow displacement.
The two positive densities then persist in one common
open interior set. Angular dependence is allowed
throughout; no symmetry removes the mutual evolution
of the dusts.

\section{The exact null current at fixed edges}
\label{sec:current}

For \(h=\delta g\) use the bare Einstein Lee--Wald
potential, with vector components
\begin{equation}\label{eq:theta}
 \theta^\mu(g;h)
 =\frac{1}{16\pi G}
   \left(\nabla_\nu h^{\mu\nu}-\nabla^\mu\tr_g h\right).
\end{equation}
The associated spacetime form contracts this vector into
the metric volume form, and the two-form on solutions is
its exterior derivative in the indicated family
\cite{LW}. We use no additional exact boundary shift
of this potential. On one initial face let \(s\) be its
fixed affine generator label and \(r\) its transverse
null coordinate. In particular the reconstructed shift
on \(H_1\) lies in \(g_{rA}\) and \(g_{rr}\), not in
\(g_{sA}\) or \(g_{ss}\). On either face,
\[
 g_{ss}=g_{sA}=0,\qquad g_{sr}=-c,\qquad
 g^{r\mu}=-c^{-1}\delta_s^\mu,\qquad c=2\Omega^2=2 .
\]
The shift and transverse metric can vary, but
\(h_{s\mu}=h^{r\mu}=0\) on the face. Affineness gives
\(\partial_r g_{ss}=0\), and a tangent of the affine
family has \(\partial_r h^{rr}=c^{-2}\partial_r h_{ss}=0\)
there. Terms from differentiating the inverse metrics in
this last identity vanish because \(h_{s\mu}=0\).
The connection terms with a lower \(s\) vanish:
\(\Gamma^r_{ss}=\Gamma^r_{sA}=0\). The remaining
component is \(\Gamma^r_{AB}=q'_{AB}/(2c)\).
Also \(-\nabla^r\tr_g h=c^{-1}\partial_s
\tr(q^{-1}\delta q)\) and the metric density is
\(c a^2\sqrt{\det\sigma}\). Thus the integrand, in
units \(1/(16\pi G)\) and relative to
\(\dd s\,\dd A_\sigma\), is
\[
 a^2\left\{\frac12\tr\big[
       (q^{-1}q')(q^{-1}\delta q)\big]
       +\partial_s\tr(q^{-1}\delta q)\right\}.
\]
Since \(q=a^2\widehat q\) and the conformal determinant
is fixed, both \(B_\lambda=\widehat q^{-1}\widehat q'\)
and \(B_\delta=\widehat q^{-1}\delta\widehat q\) are
traceless and \(\tr(q^{-1}\delta q)=4\delta a/a\).
The expansion-cross term from the product cancels the
derivative of the denominator in the last trace.

To be explicit about signs, orient the null base by
\(\dd u\wedge\dd v\) and the sphere by its marked
orientation. The pullback on \(H_1\) selects
\(+\theta^u\,\dd v\), while on \(H_2\) it selects
\(-\theta^v\,\dd u\), with the area densities understood.
A consistently oriented Cauchy arc runs in decreasing
\(u\) along \(H_2\) and increasing \(v\) along \(H_1\).
Thus, when both integrals are written in their own
increasing \(s_i\), the signs agree. We obtain the exact
null primitive and its derivative
\begin{align}
 \Theta_i
 &=\frac{1}{16\pi G}\int_{H_i}
 \left(4a_i\,\dd_{\!P}a_i'
   +\frac{a_i^2}{2}\tr(B_{\lambda,i}B_{P,i})\right)
   \dd s_i\,\dd A_{\sigma_i},
 \qquad
 B_{P,i}=\widehat q_i^{-1}\dd_{\!P}\widehat q_i ,
 \label{eq:primitive}\\
 \mathfrak S_\lambda&=\dd_{\!P}(\Theta_1+\Theta_2).
 \nonumber
\end{align}
This is a direct affine specialization of the
area and shear pairings in \cite{Reisenberger,HF},
with no field-dependent angular relabeling of the
marked sheets.

The full spacetime symplectic current is closed for
vacuum solutions and their linearized solution tangents.
A monotone base arc with \(\dot u\dot v<0\) and the same
two endpoints, fibered by the sphere, is a homologous
spacelike surface: on the orthogonal lift of the arc its
metric is \(-4\Omega^2\dot u\dot v>0\), and the angular
metric is positive. The region between it and the initial
two-face surface has no lateral face. Stokes' theorem
therefore transports \(\mathfrak S_\lambda\) to this cut
at every fixed wavelength. The joint contributes no
additional corner term in this bare-current prescription.
The outer embeddings are fixed, not all metric or
transverse corner jets. Diffeomorphisms vanishing in
neighborhoods of the retained outer edges are null
directions by the boundary Noether identity. We do not
use such a degeneracy for diffeomorphisms which move or
act nontrivially at the edges. Differently cut surfaces
require their actual lateral flux and any specified
corner prescription.

\section{Limit and the data retained by it}\label{sec:limit}

It is the derivative of the area one-form that is
negligible: in \eqref{eq:primitive} its exterior derivative
is \(4\dd_{\!P}a\wedge\dd_{\!P}a'=O(\lambda^3)\) by
the area lemma. The primitive itself need not be
\(o(\lambda)\). For the shape block,
\[
 B_P=\lambda F\,\re((\dd_{\!P}Z)e^{is/\lambda})
          +O(\lambda^2),\qquad
 B_\lambda=A+F\re(iZe^{is/\lambda})+O(\lambda).
\]
The remainders have the required differentiated bounds
on the selected profile families. Use
\(a_\lambda=\bar a+O(\lambda^2)\).
The product involving \(A\) is a first harmonic whose
integral, after division by \(\lambda\), is \(O(\lambda)\).
The two fast factors have common zero harmonic
\begin{equation}\label{eq:shape-mean}
 {\rm av}_{\vartheta}\,
 \tr\!\left[\re(iZ e^{i\vartheta})
            \re((\dd_{\!P}Z)e^{i\vartheta})\right]
   =\tfrac12\,\im\tr(Z^\dagger\dd_{\!P}Z)
   =\tfrac12\,\im(z^\dagger\dd_{\!P}z).
\end{equation}
The other harmonics integrate to \(O(\lambda)\) after
scaling. These estimates hold also after one profile
derivative of the scaled shape primitive, so exterior
differentiation gives \eqref{eq:theorem}. In particular
the unscaled two-form vanishes in the limit on these
bounded tangents.

On a regular projective chart,
\(\im(z^\dagger\dd_{\!P}z)=\dd_{\!P}\chi+
p\dd_{\!P}\phi\); its differential is
\(\dd_{\!P}p\wedge\dd_{\!P}\phi\), the global curvature
of the unit-vector phase bundle. Choose the ordered
tangents with \(\delta_1p=\psi,\delta_2\phi=\psi\)
and other coordinates fixed on a patch where
\(0<p<1\) and \(F\ne0\). The pairing is a positive
constant times the integral of \(\psi^2\) against
the positive fixed intensity, hence nonzero for
nonzero compactly supported \(\psi\). A tangent
on one sheet and one on the other have zero
cross pairing already in the exact pullback
\eqref{eq:primitive}. Each primitive depends only
on the intrinsic geometry of its own face and the
fixed marks. It does not follow that the two streams
have decoupled vacuum or dust dynamics in the interior.

The positive-generator-frequency polarization matrix
on each face is proportional to \(zz^\dagger\); its
normalized rank-one part fixes the ray \([z]\). The
density of \eqref{eq:theorem} is precisely
\((8\pi G)^{-1}\nu_i^{\rm in}\) multiplying the
projective form. In this sector, scalar dust therefore
supplies the weight but not the variables that carry
the pairing. The exact solution family has invisible
directions for the scalar-dust map and a nonzero
normalized canonical pairing between selected such
directions. Keeping the projective field on each
incident sheet restores this pairing. The local
overall carrier phases are radical directions of the
normalized restriction; no finite-wavelength pure
gauge conclusion is drawn from this fact.

This theorem is on marked, fixed-edge cuts and a chosen
coherent single-carrier sector. It neither constructs a
pointwise Jones field of a general spacelike weak limit
nor asserts a multi-eikonal symplectic limit, a theory of
arbitrary moving edges, or a quantum-state construction.

\begin{thebibliography}{9}

\bibitem{LR}
J.~Luk and I.~Rodnianski,
\emph{High-frequency limits and null dust shell solutions
in general relativity}, \href{https://arxiv.org/abs/2009.08968}{arXiv:2009.08968}.

\bibitem{LW}
J.~Lee and R.~M.~Wald,
\emph{Local symmetries and constraints},
J.\ Math.\ Phys.\ \textbf{31}, 725--743 (1990),
\href{https://doi.org/10.1063/1.528801}{doi:10.1063/1.528801}.

\bibitem{Reisenberger}
M.~P.~Reisenberger,
\emph{The symplectic 2-form for gravity in terms of free
null initial data}, \href{https://arxiv.org/abs/1211.3880}{arXiv:1211.3880}.

\bibitem{HF}
F.~Hopfm\"uller and L.~Freidel,
\emph{Gravity degrees of freedom on a null surface},
\href{https://arxiv.org/abs/1611.03096}{arXiv:1611.03096}.

\bibitem{AS}
A.~Ashtekar and T.~A.~Schilling,
\emph{Geometrical formulation of quantum mechanics},
\href{https://arxiv.org/abs/gr-qc/9706069}{arXiv:gr-qc/9706069}.

\bibitem{paper29}
\emph{Polarization transport in the high-frequency limit of the vacuum Einstein equations},
\href{https://github.com/g3173028019-beep/ai-automated-quantum-gravity-research/tree/main/papers/29-high-frequency-graviton-polarization}{companion Paper 29}.

\bibitem{paper29B}
\emph{Physical polarization of high-frequency vacuum gravity: transport, exact preparation, and curvature readout},
\href{https://github.com/g3173028019-beep/ai-automated-quantum-gravity-research/tree/main/papers/29b-high-frequency-graviton-polarization}{companion Paper 29B}.

\bibitem{transfer}
\emph{Polarization defects and opposite graviton transfers in exact vacuum sandwiches with identical geodesic returns}, edition 2,
\href{https://github.com/g3173028019-beep/ai-automated-quantum-gravity-research/tree/main/papers/29b-31-polarization-defects-graviton-transfer}{companion Paper 29B+31+Q19}.

\end{thebibliography}
\end{document}
